\Needspace{8\baselineskip}
\Article{271}
It is now clear that all ternary forms, just like binary forms, can be divided into \textit{classes}, assigning equivalent forms to the same class and inequivalent forms to different classes. Forms with different determinants will therefore certainly belong to different classes, so there will be infinitely many classes of ternary forms. Ternary forms with the same determinant, however, make up sometimes fewer and sometimes more classes. But the outstanding property of these forms is that \textit{all forms with the same given determinant always make up a f\,i\,n\,i\,t\,e number of classes.} Before developing this very important theorem more fully, we must describe the following essential distinction among ternary forms.

Some ternary forms are such that both positive and negative numbers can be represented by them without distinction, \textit{for example} the form $xx+yy-zz$. They will therefore be called \textit{indefinite forms}. By others, on the contrary, negative numbers cannot be represented, but only positive ones (apart from zero, obtained by setting each indeterminate $=0$), as with $xx+yy+zz$, so they will be called \textit{positive forms}. Finally, by others positive numbers cannot be represented, as with $-xx-yy-zz$, so they will be called \textit{negative forms}. Positive and negative forms will be given the common name \textit{definite forms}. Here are general criteria by which these properties of forms can be distinguished.

Multiplying the ternary form
\[
f=axx+a'x'x'+a''x''x''+2bx'x''+2b'xx''+2b''xx'
\]
\SourceLeaf{311}
of determinant $D$ by $a$, and denoting the coefficients of the adjoint of $f$, as in Article 267, by $A$, $A'$, $A''$, $B$, $B'$, $B''$, gives
\[
(ax+b''x'+b'x'')^2-A''x'x'+2Bx'x''-A'x''x''=g
\]
Multiplying again by $A'$ gives
\[
A'(ax+b''x'+b'x'')^2-(A'x''-Bx')^2+aDx'x'=h
\]
From this we immediately conclude that if both $A'$ and $aD$ are negative numbers, all values of $h$ are negative. Clearly, then, the form $f$ can represent only numbers whose sign is opposite to that of $aA'$, \textit{that is}, the same as that of $a$, or opposite to that of $D$. In this case, therefore, $f$ will be a definite form, positive or negative according as $a$ is positive or negative, or according as $D$ is negative or positive.

If, however, either both $aD$, $A'$ are positive, or one is positive and the other negative (neither $=0$), it is easy to see that $h$ can take both positive and negative values by a suitable choice of the quantities $x$, $x'$, $x''$. In this case, then, $f$ can take values with both the same sign as $aA'$ and the opposite sign, and so it will be an indefinite form.

In the case where $A'=0$, but not $a=0$, we have
\[
g=(ax+b''x'+b'x'')^2-x'(A''x'-2Bx'')
\]
Giving $x'$ an arbitrary value (but not $=0$), and taking $x''$ so that $\frac{A''x'}{2B}-x''$ has the same sign as $Bx'$ (which is easily seen to be possible, since $B$ cannot be $=0$, as this would give $BB-A'A''=aD=0$, and therefore also $D=0$, a case we exclude), makes $x'(A''x'-2Bx'')$ a positive quantity. It is then easy to see that $x$ can be chosen so that $g$ takes a negative value. Clearly these values can also be chosen so that, if desired, they are all integers. Finally, it is clear that whatever values are assigned to $x'$, $x''$, $x$ can be taken so large that $g$ becomes positive. We conclude that in this case the form $f$ is indefinite.
\SourceLeaf{312}
Finally, if \(a=0\), we will have
\[
f=a'x'x'+2bx'x''+a''x''x''+2x(b''x'+b'x'')
\]
Taking \(x'\), \(x''\) arbitrarily, but so that \(b''x'+b'x''\) is not \(=0\) (which is clearly possible unless both \(b'\) and \(b''\) are \(=0\); but then we would have \(D=0\)), we readily see that \(x\) can be chosen so that \(f\) takes both positive and negative values. So in this case too \(f\) will be an indefinite form.

Just as here we have used the numbers \(aD\), \(A'\) to judge the nature of the form \(f\), we can also use \(aD\) and \(A''\): \(f\) is a definite form if both \(aD\) and \(A''\) are negative, and indefinite in every other case. In exactly the same way, considering the numbers \(a'D\) and \(A\), or \(a'D\) and \(A''\), or \(a''D\) and \(A\), or finally \(a''D\) and \(A'\), can serve the same purpose.

From all this it follows that in a definite form the six numbers \(A\), \(A'\), \(A''\), \(aD\), \(a'D\), \(a''D\) are negative. In a positive form, \(a\), \(a'\), \(a''\) will be positive and \(D\) negative; in a negative form, \(a\), \(a'\), \(a''\) will be negative and \(D\) positive. Thus all ternary forms with a given positive determinant fall into negative and indefinite forms, and all those with a negative determinant into positive and indefinite forms. Finally, there are no positive forms with a positive determinant, or negative forms with a negative determinant. --- From the same considerations it is easy to see that the adjoint of a definite form is always definite, and indeed \emph{negative}, while the adjoint of an indefinite form is indefinite.

Since all numbers representable by a given ternary form can clearly also be represented by every form equivalent to it, the ternary forms contained in the same class will either all be indefinite, all positive, or all negative. We may therefore extend these names of forms to entire classes as well.
